% Exploratory AI-agent evidence; do not describe as a human study.
\subsection{Exploratory Multi-Model Agent Replication}
\label{sec:gateway-agent-replication}

We examined whether OpenCoder's uncertainty presentation also changes the
behavior of automated code reviewers. Twelve gateway-mediated model
configurations each reviewed the same 12 frozen ExecRepoBench outputs under a
counterbalanced design, yielding 144 planned episodes (48 per condition). The
models received no tools, test outcomes, or conversational history, and every
returned function was evaluated with the same private executable tests. This
analysis is exploratory: gateway models are neither human participants nor
equivalent to the corresponding consumer chat products.

Across the final campaign, 134/144
episodes produced executable outputs. Generic review achieved
38.6\% correctness among evaluable
outputs, compared with 46.7\% for the
uncertainty display and 35.6\% for
targeted guidance. When model-side non-response is retained in the planned
denominator, the corresponding end-to-end success rates were
35.4\%,
43.8\%, and
33.3\%. Relative to generic review,
targeted guidance changed end-to-end success by
-2.1 percentage points (model-clustered
95\% bootstrap CI [-12.5,
+8.3]); the uncertainty-only
display changed it by +8.3 points
([-2.1,
+18.8]). These intervals are
descriptive sensitivity estimates over the twelve configured systems, not
population-level inference over all LLMs.

\noindent\textbf{Finding.}
The multi-model replication does not show a positive aggregate benefit from source-specific guidance. The result supports treating
AI-agent evidence as a robustness check while keeping human usability and
developer effectiveness as separate empirical questions.
